\cotefile{84}
\scanpage{19}

\[
\boxed{\delta(u)\left(\mathrel{\overset{\mathrm{déf\unsure{n}}}{=}}(\operatorname{Tr}u)^2-4N(u)\right)=-\underbrace{\delta(Q_u)}_{\substack{\text{discr. de la}\\\text{forme }Q_u}}}
\]

-- Attention au signe $-$ dans le second mb\unsure{!|!.}

On note que si $Q$ est une forme quadratique à valeurs sur le module \deleted{inv. \unsure{$\Delta$} $\Lambda$, alors}\inserted{inv. $\deleted{M}\Delta'$ \deleted{(voir not. \unsure{inv.})}}, alors \deleted{\unsure{$\varphi$}}

\[
\delta(Q)\in\Gamma\left({\deleted{M}\Delta'}^{\otimes2}\otimes L^{\otimes2}\right)
\]

est défini ainsi :

\[
\varphi_Q:E\otimes E\longrightarrow\deleted{\unsure{M}}\Delta'
\]

d'où

\[
E\longrightarrow \underline{\operatorname{Hom}}(E,M)\simeq E^\vee\otimes\deleted{M}\Delta',
\]

\deleted{\unsure{et}}d'où en prenant l'ext. \unsure{ème|âme|deux} \deleted{\unsure{le} dét} module dét $=\wedge^2$ :

\[
\underbrace{\det E}_{L}\mathbin{\unsure{-|\longrightarrow}}\underbrace{\det(E^\vee)}_{L^\vee}\otimes{\Delta'}^{\otimes2}
\]

d'où,

\[
\delta(Q)\in\Gamma\,\underbrace{L^{\vee\otimes2}{L'}^{\otimes2}}_{(\Delta'\otimes L^\vee)^{\otimes2}}\deleted{\unsure{-}}.
\]

\deleted{Mais d'autre part on a} Quand $\Delta'\simeq L$, c'est une section de $\underline{\mathcal O}$, ce qui donne un sens à la formule sur les discriminants \ldots

2.4. Soit

\[
Q:E\mathbin{\unsure{-|\longrightarrow}}\Delta'\qquad Q\in\Gamma\,\underline{\operatorname{Quad}}(E,\Delta')
\]
% archive sheet number: 18
\scanpage{20}

une forme quadratique à valeurs dans le module inversible $\Delta'\inserted{=\Delta^\vee}$. Il revient au $=$ de se donner

\[
\Delta\longrightarrow\underline{\operatorname{Quad}}(E)\qquad
\left(\simeq\underline{\operatorname{End}}(E)/\underline O\cdot1\unsure{)}\deleted{\unsure{\otimes}}\otimes L^\vee\right)
\]

(vrai pour \unsure{tt} $\Delta$ loc. libre de r. f.)

i.e. de se donner

\[
u_Q:\underbrace{\deleted{\unsure{M}}\Delta\otimes L}\longrightarrow\underline{\operatorname{End}}(E)/\underline O\cdot1
\]

Nous allons en profiter pour définir une alg. \deleted{\unsure{pr.}} quadratique $\mathcal A_Q$ dans $X$, et un hom \inserted{d'alg.}

\[
\mathcal A_Q\mathbin{\unsure{-|\longrightarrow}}\underline{\operatorname{End}}(E)
\]

dont l'image soit l'anneau $\mathcal A_\Sigma$ des similitudes pour la structure de similitude $\Sigma\inserted{=\Sigma_Q}$ définie par $Q$, \inserted{i.e l'image inv. de $\operatorname{Im}u_Q$ dans $\underline{\operatorname{End}}(E)$.} Pour ceci, on \deleted{\unsure{regarde}} la structure d'extension

\[
0\longrightarrow\underline O_X\longrightarrow\underline{\operatorname{End}}(E)\longrightarrow\underline{\operatorname{End}}(E)/\underline O\cdot1_E\longrightarrow0
\]

et \inserted{on désigne par $\mathcal A_Q$} l'image inverse de cette extension par $u_Q$, donc on a un hom d'extensions

\diagram{\begin{tikzcd}[ampersand replacement=\&,column sep=small]
0\arrow[r,no head]\&\underline O\arrow[r,no head]\arrow[d,"{\mathrm{id}}"]\&\mathcal A_Q\arrow[r,no head]\arrow[d,"{u'_Q}"]\&\overset{\Delta\otimes L}{\overset{\scriptscriptstyle\Vert\,\mathrm{déf\unsure{n}}}{\Lambda}}\arrow[r,no head]\arrow[d,"{u_Q}"]\&0\\
0\arrow[r,no head]\&\underline O\arrow[r]\&\underline{\operatorname{End}}(E)\arrow[r,no head]\&\operatorname{End}(E)/\underline O\arrow[r]\&0
\end{tikzcd}}
% archive sheet number: 19
\scanpage{21}

Considérons sur $\mathcal A_Q$ la forme quadratique $N_Q$ définie par \deleted{\unsure{la}}

\[
N_Q(\deleted{\unsure{u}}x)=N(u'_Q(x))
\]

on aura donc, en désignant par $1_Q$ la section $1$ de $\underline O_{\unsure{\alpha|X}}$, regardée comme section de $\mathcal A_Q$

\[
N_Q(1_Q)=1,
\]

Ainsi,

\[
\deleted{\unsure{A}}(\mathcal A_Q,1_Q,N_Q)
\]

est un module loc. libre de rang 2, muni d'une section $1_Q$ et d'une forme quadratique valant $1$ sur cette section. On a donc une structure d'alg. quadratique unique sur $\mathcal A_Q$, pour laquelle $N_Q$ soit la norme. De plus, on vérifie immédiatement que $u'_Q$ est un hom d'algèbres, et qu'on a, en plus de la relation sur les normes, la relation

\[
\begin{aligned}
\deleted{\unsure{N}}\operatorname{Tr}(u'_Q(x))&=\operatorname{Tr}_{\mathcal A_Q/\underline O}(x)\\
u'_Q(\bar x)&=u'_Q(x)^{\text{--}}
\end{aligned}
\]

\unsure{ici|où} on désigne par

\[
v\longmapsto\bar v=(\operatorname{Tr}v)\cdot1-v
\]

l'antiinvolution canonique de $\underline{\operatorname{End}}(E)$, égale à $\mathrm{id}_{\underline O}$ sur $\underline O\cdot1_E$ et induisant $-\mathrm{id}$ sur le quotient.
% archive sheet number: 20
\scanpage{22}

\uuline{NB} $1^\circ$. La corr. $(E,\Delta,Q)\longmapsto\deleted{\{}\mathcal A_Q\deleted{\unsure{\mathcal A_Q\subset|\mathcal A_Q\to}\underline{\operatorname{End}}(E)}$ est \uline{fonctorielle}\inserted{(pour les iso)}, et l'hom (pas néc. inj.)

\[
\mathcal A_Q\longrightarrow\underline{\operatorname{End}}(E)
\]

\deleted{est} fonctoriel. De plus, ces constructions commu\deleted{\unsure{tent}} aux changements de base.

$2^\circ$. Il en résulte, par transport de structure, que \inserted{pour $\lambda\in\Gamma\,\underline O_X^*$ (scalaires \uline{inversibles})}, on a un isom can. d'alg. quadratiques

\[
\mathcal A_Q\simeq\mathcal A_{\lambda Q}
\]

compatible avec les actions sur $E$. Bien sur, il n'en est pas ainsi si on cesse de supposer $\lambda$ inversible (prendre $\lambda=\deleted{\unsure{1}}0$ !).

$3^\circ$) Conditions équivalentes :

a) $\mathcal A_Q\longrightarrow\underline{\operatorname{End}}(E)$ univ. inj. (i.e. inj. sur chaque fibre)

b) $Q$ n'est nulle sur aucune fibre réduite.

c) $\Delta\overset{Q}{\longrightarrow}\operatorname{Quad}(E)$ est univ. injectif

d) $\Sigma_Q$ est strictement fidèle.

(C'est essentiellement le lemme 2 plus haut.) Alors la structure d'Alg. sur $\mathcal A_Q$ est simplement celle induite par $\underline{\operatorname{End}}(E)$.

2.5. Cas tot. dégénéré de la forme \uline{nulle} $Q:E\mathbin{\unsure{-|\longrightarrow}}\deleted{\unsure{L}}\Delta'$. Alors

\[
\mathcal A_Q=\underline O\times\Lambda\qquad(\Lambda=\Delta\otimes L)
\]

avec la structure d'algèbre définie par la condition que $\Lambda$ est de \uline{carré nul}, i.e.
% archive sheet number: 21
\scanpage{23}

\[
\mathcal A_Q\simeq\operatorname{\unsure{Duaux|Deux}}_{\underline O}(\Lambda).
\]

L'hom

\[
u'_Q:\mathcal A_Q\mathbin{\unsure{-|\longrightarrow}}\underline{\operatorname{End}}(E)
\]

est \deleted{défini}\inserted{déterminé} par la condition \deleted{que} de s'annuler sur $\Lambda$, i.e. c'est $(a,\lambda)\longmapsto a$. \deleted{\unsure{Elle}} Il n'a rien d'injectif.

2.6. Quand on se donne une base de $E$ et de $\Delta$, alors $Q$ est donnée par

\[
Q(\lambda,\mu)=\alpha\lambda^{\unsure{1|2}}+\beta\lambda\mu+\gamma\mu^2
\qquad\alpha,\beta,\gamma\in\Gamma\underline O,
\]

et on \deleted{peut prendre} a une \deleted{\unsure{syst.|base}}\unsure{base} $(1,u)$ de $\mathcal A_Q$, avec

\[
\left|\begin{aligned}
u^2-\beta u+\alpha\gamma&=0\qquad\text{i.e. }\deleted{\beta=\operatorname{Tr}u},\\
\operatorname{Tr}u&=\beta\\
Nu&=\alpha\gamma
\end{aligned}\right.
\]

L'image de $u$ dans $\operatorname{End}(E)\simeq\mathbb M_2(\underline O)$ est la matrice déjà envisagée

\[
u_Q=\begin{pmatrix}0&-\gamma\\\alpha&\beta\end{pmatrix}
\]

Pour achever de déterminer $\mathcal A_Q$ comme foncteur en $(E,\Delta,Q)$, il faut voir comment s'explicite l'isom de $A_{\alpha,\beta,\gamma}=A[u]/(u^2-\beta u+\alpha\gamma)$ avec $A_{\alpha',\beta',\gamma'}=A[u]/(u^2-\beta'u+\alpha'\gamma')$, quand on passe d'un syst. de bases de $(A,\Delta)$ à un autre, par des matrices $\begin{pmatrix}u_{11}&u_{12}\\u_{21}&u_{\unsure{12|22}}\end{pmatrix}$ et $v$ \unsure{..|\ldots}
% archive sheet number: 22
\scanpage{24}

Je ne vais pas faire l'explicitation ici.

2.6. Le rev. quadratique \inserted{de $X$} associé à $\mathcal A_Q$ mérite de s'appeler l'\uline{invariant d'Arf} de la forme quadratique binaire $Q$. On définira un inv. d'Arf pour \uline{\unsure{tte|les} forme} quadratique scalaire sur un module loc. libre de rang donné (mais en distinguant les cas du rang pair et du rang impair, pour la nature de cet invariant \ldots).

3. \uline{Théorème récapitulatif}

Equivalences entre \deleted{\unsure{deux}} catégories fibrées sur \inserted{la catégorie} $\mathrm{Sch}$\deleted{\unsure{ém}} des schémas $X$ :

\textcircled{1} Structures de similitudes binaires \inserted{str. fidèles} sur $X$, i.e. des systèmes

\[
(E,\deleted{\unsure{Q}}\Delta),
\]

où $E$ est un \deleted{\unsure{M}}module loc. libre de rang $2$ sur $X$, et

\[
\Delta\subset\underline{\operatorname{Quad}}(E)
\]

un sous-\deleted{Module \unsure{inver.}}\inserted{fibré vectoriel} \uline{inversible} de $\underline{\operatorname{Quad}}(E)$, i.e. un sous-module inversible \inserted{et} loc. fact. direct (ou encore, tel que $\Delta\longrightarrow\underline{\operatorname{Quad}}\,E$ soit univ. injectif \unsure{--}).
% archive sheet number: 23
\scanpage{25}

\textcircled{$1'$} Systèmes

\[
(E,\Delta',Q),
\]

où $E$ est comme dans $1^\circ$, $\Delta'$ un Module inv,

\[
Q:E\mathbin{\unsure{-|\longrightarrow}}\Delta'
\]

une forme quadratique, telle que $Q$ soit non \inserted{(identiquement)} nulle sur tte fibre réduite.

\textcircled{2} Systèmes $\deleted{\unsure{(X',\Lambda')}}\,(\mathcal A,E)$ d'une alg. quadratique $\mathcal A$ sur $\underline O$, et d'un $\mathcal A$-module inversible i.e. loc. libre de rang $1$.

\textcircled{$2'$} Systèmes $(X',\Lambda')$ d'un rev. quadratique $X'$ de $X$, et d'un Module inversible $\Lambda'$ sur $X'$.

\uline{NB} $1\Longleftrightarrow1'$ trivial, $2\Longleftrightarrow2'$ est \unsure{standart|standard} par dictionnaire des spectres d'Anneaux \unsure{qu. coh.|quasi-coh.} $(1,1')\Longrightarrow2$ a été vu plus haut.

Il faut expliciter \textcircled{$2'$} $\longrightarrow$ \textcircled{$1'$}. On prendra

\[
\left\{\begin{aligned}
 E&=p_*(\Lambda'),\quad\text{où }p:X'\longrightarrow X\quad\text{\unsure{morph. str.|morph. struct.}}\\
 \Delta'&=N_{X'/X}\,\Lambda'\\
 Q&\mathrel{\deleted{=}:}p_*(\Lambda')\longrightarrow\Delta'=N_{X'/\unsure{X|x}}(\Lambda')
\end{aligned}\right.
\]

l'\deleted{\unsure{appl}}application norme
% archive sheet number: 24
\scanpage{26}

\uline{Cor. 1} Ces catégories équivalent aussi à celle \inserted{\textcircled{$3^\circ$}} des systèmes

\[
(\mathcal A,P),
\]

où $\mathcal A$ est une alg. quadratique sur $X$, et où $P$ est un torseur\textsuperscript{*} sous le schéma en groupes comm. lisse \inserted{de dim relative $2$} sur $S$

% slanted
\margin{* torseur pour \unsure{fppf|fpqc} ou pour (Zar), \unsure{car|et} cela revient ici au $=$, \unsure{ou|où} pour tte top. intermédiaire sur $\mathrm{Sch}/S$}

\[
W(\mathcal A)^*=\left\{S_{/X}\longmapsto\Gamma(S,\mathcal A_S^*)\right\}
\]

(où $\mathcal A_S$ désigne l'image inverse de $\mathcal A$ sur $S$).

\deleted{En termes de}\inserted{En termes} de $(E,Q)$, $P$ est le schéma sur $X$ des sections locales de $E$, telles que $Q(x)$ soit une base de $\Delta'$.

D'autre part, $\Delta'$ s'explicite en termes de $(\mathcal A,P)$ comme le fibré en droites dont le fibré principal associé est

\[
P\mathbin{\unsure{\wedge|\times}}^{\unsure{W(\mathcal A)^*}}G_m,
\]

via l'hom norme

\[
W(\mathcal A)^*\xrightarrow{N_{\mathcal A/\underline O}}G_m,
\]

\unsure{s'inscrivant|s'insérant|s'insère} dans la suite exacte

\[
\unsure{1|0}\longrightarrow\underset{\substack{\text{schéma en}\\\text{groupes }\text{``\unsure{Unitaire|unitaire}''}\text{ des}\\\text{éléments de norme }1\\\text{de }\mathcal A}}{U_{\mathcal A}}\longrightarrow W(\mathcal A)^{\unsure{\times|*}}\xrightarrow{N}G_m\longrightarrow1
\]
% archive sheet number: 25
\scanpage{27}

\uline{Cor. 2} \unsure{Èquiv.|Équiv.} de cat. fibrées entre

$1^\circ$) $(E,Q)$, $Q$ une forme quadratique binaire qui n'est nulle sur aucune fibre réduite

$2^\circ$) Systèmes $(\mathcal A,\Pi)$, où $\mathcal A$ est une alg. quadratique sur $S$, et où $\Pi$ est un torseur\textsuperscript{*} sous le groupe unitaire $U_{\mathcal A}$. \deleted{\unsure{*}}

% slanted
\margin{Torseur \unsure{sens|sous} fppf ou \unsure{sens|sous} (ét), cela revient ici au $=$, mais la top. (Zar) n'est plus assez fine \ldots}

$\Pi$ est le sous-schéma de $E^{\unsure{*}}$ des sections locales $x$ telles que $Q(x)=1$, qui \deleted{sont}\inserted{forment bien} un torseur sous $U_{\mathcal A}$.

\uline{NB} Le rev. quadratique associé à la forme quadratique $Q:E\longrightarrow\Delta'$ (non nulle sur toute fibre) peut s'expliciter comme \inserted{le} sous-schéma de ${\unsure{P|\mathbb P}}^1(E)$ défini par l'équation $Q(x)=0$. Si $E$ est considéré comme Module inv. sur l'Algèbre quadratique $\mathcal A$, le morphisme

\[
\operatorname{Spec}(\mathcal A)\lhook\joinrel\longrightarrow{P}^{\unsure{*1|1}}(E)
\]

correspondant \deleted{peut} s'explicite ainsi : pour toute section locale $u$ de $\operatorname{Spec}(\mathcal A)$ sur $A$ i.e.
% archive sheet number: 26
\scanpage{28}

pour t\deleted{\unsure{out}} hom d'algèbres $\mathcal A\longrightarrow\underline O$, le noyau $J\subset\mathcal A$ est un sous-fibré vectoriel \inserted{de rang $1$ \deleted{\unsure{(suppl. de $\underline O$ $\cdot1_{\mathcal A}$)}}} sur lequel $N_{\mathcal A/O}$ est nul, \inserted{(et inversement)} et $D=J\cdot E$ est alors un sous-fibré en droites de $E$ tel que $Q\mid D=0$. Inversement tt sous-fibré en droites de $E$ ayant cette propriété provient d'une section unique de $X'$ sur $X$. (Pour le vérifier, OPS $E=\mathcal A$, et $Q$ la norme, et c'est alors immédiat.)

4. La notion de structure de similitude sur $E$, et la distinction de celles qui sont fidèles, str. fidèles, \inserted{voire} non dégénérées, ``ne \unsure{dépend|dépendent} que de $E$ modulo tensorisation par un module inv'', donc ne \unsure{dépend|dépendent} que de l'algèbre d'Azumaya $\mathcal M$ associée,

\[
\mathcal M=\underline{\operatorname{End}}(E),
\]

de rang $4$ sur $E$. Si on part de une telle $\mathcal M$, la donnée d'une structure de similitude \deleted{relative à un Module \unsure{+}} équivaut à celle d'un sous-Module loc. \unsure{mon.|monog.|monogène} \inserted{$\Delta$} \unsure{de|du}

\[
\mathcal M/\underline O\,1_{\mathcal M}\overset{\mathrm{déf\unsure{n}}}{=}\Sigma(\mathcal M),
\]

(qui joue le rôle du Module des formes quadratiques)
% archive sheet number: 27
\scanpage{29}

Le cas \deleted{\unsure{fidèle}} fidèle (resp. str. fidèle) est celui où ce sous-Module est un Module inversible, resp. un sous-fibré en droites. Dans ce dernier cas, il \deleted{\unsure{y}} y a un \deleted{fibré}\inserted{rev.} quadratique associé de $X$, qui peut s'interpréter comme l'intersection de $\unsure{\check P|\check{\mathbb P}}(\mathcal A)\subset\check P(\mathcal M)$ (où $\mathcal A\inserted{\supset\underline O\cdot1_{\mathcal M}}$ est le sous-Anneau de $\mathcal M$ image inverse de $\Delta$) avec la quadrique

\[
\mathcal Q\subset\check P(\mathcal M)
\]

définie par l'équation $\operatorname{Nr}_{\mathcal M/\unsure{k|h}}(x)=0$, où $\operatorname{Nr}$ désigne la norme réduite. (C'est tautologique, vu que ladite norme induit sur $\mathcal A$ celle de $\mathcal A$.)

Que devient le $\mathcal A$-Module inv \inserted{$E$}, i.e. le $\underline O_{X'}$-Module inversible $\Lambda$, de \unsure{la théorie précédente|le théorème précédent} ?

\unsure{Il faut|Il} s'attendre à ce qu'il ne soit plus défini que modulo remplacement de $E$ par $E\otimes_{\underline O_X}\Theta$, $\Theta$ inv. sur $X$, i.e. que $\Lambda$ soit défini mod changement en $\Lambda\otimes p^*(\Theta)$, où $p:X'\mathbin{\unsure{-|\longrightarrow}}X$ le morphisme structural. Mais si $\sigma$ désigne l'inv. can. de $X'/X$, alors $\Lambda\sigma{\unsure{(\Lambda)|({}^*\Lambda)|^*(\Lambda)}}^{-1}$ reste invariante, par changement de $\Lambda$ en $\Lambda\otimes p^*(\Theta)$. D'autre part,
% archive sheet number: 28
\scanpage{30}

On a

\[
N(\Lambda\sigma(\Lambda)^{-1})\simeq N(\Lambda)(N(\sigma(\unsure{A|\Lambda})))^{-1}\simeq\underline O,
\]

isom can dû au fait que

\[
N(\sigma(\unsure{A|\Lambda}))\simeq N(\Lambda)\qquad\text{canoniquement.}
\]

Donc on s'attend à trouver un $A$-Module loc libre, avec restriction du groupe structural de $W(A)^*$ à $U(A)$, \inserted{schéma en} groupes des sections locales de $A$ de norme $1$.

Le $\mathcal A$-Module qui se propose à la vue est évidemment $\mathcal M/\mathcal A$, lequel est bien libre de rang $1$ sur $\mathcal A$, si $\mathcal A$ est un sous-fibré vect. de rang $2$ \deleted{de $\mathcal M$} (cas ``str. fidèle''). \deleted{Mais}

Dans le cas où $\mathcal M=\underline{\operatorname{End}}(E)$, il faudrait voir qu'on a un ``iso. à scalaire près'' canonique

\[
E\xrightarrow{\sim}\underline{\operatorname{End}}(E)/\mathcal A
\]

(iso. de $\mathcal A$-modules, mais le ``scalaire près'' réfère à une section de $\underline O_X^*$ \ldots), \deleted{ou encore un iso de schémas}

\[
\deleted{{\unsure{P|\mathbb P}}^{\unsure{\vee|*}}(E)\simeq P\bigl(\underline{\operatorname{End}}(E)^{\unsure{*|+}}/\mathcal A\bigr)}
\]

% left margin, vertical strokes
\margin{$\begin{gathered}|\\|\\|\\\unsure{!|\vert}\\)\end{gathered}$}

% slanted
\margin{\deleted{où $\underline{\operatorname{End}}^{+}$ désigne le sous-Module des end. de \uline{trace nulle}}}

\deleted{ou encore (ou mieux}

\[
\begin{aligned}
&\deleted{P(E)\simeq\unsure{\mathcal Q}\{\text{schéma des droites vect. }D\text{ dans }\underline{\operatorname{End}}(E)^+\text{ telles que }N\mid D=0\}},\\
&\deleted{D\longmapsto D^{\perp}\otimes D}
\end{aligned}
\]
% archive sheet number: 29
\scanpage{31}

il s'agit de décrire un iso. can.

\[
P(\underline{\operatorname{End}}(E)/\mathcal A)\simeq\deleted{\unsure{Q_{\mathrm{red}}}}\inserted{C,\ \text{conique can. dans }P(\underline{\operatorname{End}}^+(\mathcal M))}
\]

(\uline{NB} en termes de $E$ cependant, non en termes de $\mathcal A\subset\mathcal M$ seulement, car $\deleted{\unsure{Q}}\inserted{C}$ ne sera en général \unsure{pas} isom à un $P$(fibré vect. de rang $2$) \ldots)

Or on a sur $\underline{\operatorname{End}}(E)$ \uline{deux} structures de $\mathcal A$-module, la gauche et la droite, échangées par l'antiinvolution canonique

\[
u\longrightarrow\operatorname{Tr}(u)-u\qquad(\operatorname{Tr}=\text{trace réduite})
\]

et on a donc égalité des modules déterminants de $\underline{\operatorname{End}}(\deleted{\unsure{A}}E)$ sur \deleted{\unsure{a}} $\mathcal A$, lesquels sont égaux can. à $\underline{\operatorname{End}}(E)/\mathcal A$

\[
\underline{\operatorname{End}}(E)/\mathcal A\simeq\det_{\mathcal A}\underline{\operatorname{End}}(E)
\]

Mais l'iso

\[
\underline{\operatorname{End}}(E)\simeq E\otimes E^\vee
\]

transforme la structure de $\mathcal A$-module gauche en celle définie par l'opération $u_E\otimes\mathrm{id}_{E^\vee}$ de $u$, et la droite en $\mathrm{id}_E\otimes{}^tu_E$, \unsure{On|Donc} a \unsure{donc|ainsi|dans} \unsure{des iso|des isos|le cas} de $A$-modules

\[
\det_{\mathcal A}\operatorname{End}(E)\simeq E^{\otimes_A2}\otimes\underbrace{\bigl(\det_{\unsure{k|h}}E^\vee\bigr)}_{L^\vee}^{2}
\]

et de $=$

\[
\det_{\mathcal A}\operatorname{End}(E)\simeq\underbrace{\det_k(E)^{2}}_{L}\otimes {E^\vee}^{\otimes_A2},
\]

d'où un iso de $A$-modules

\[
\left[\operatorname{End}(E)/\mathcal A\simeq E^{\otimes_A2}\otimes_k L^\vee\simeq {E^\vee}^{\otimes_A2}\otimes_k L\right.
\]
% archive sheet number: 30
\scanpage{32}

On devine que l'iso

\[
{E^{\deleted{\unsure{\vee}}}}^{\otimes_A2}\otimes_{\unsure{k|h}}L^\vee\simeq {E^\vee}^{\otimes_A2}\otimes_k L
\]

i.e.

\[
{E^\vee}^{\otimes_A2}\simeq E^{\otimes_A2}\otimes_k{L^{\deleted{\unsure{\vee}}}}^{\otimes_k(-2)}
\]

est déduit de l'iso can. (qui respecte les structures de $A$-modules)

\[
E^\vee\simeq E\otimes_k L^{\otimes_k(-1)}
\]

en \unsure{prenant|élevant} les \deleted{\unsure{puis|par}} carrés tensoriels sur $A$.

L'iso de $A$-modules cherché (mod scalaires provenant de $\underline O_{\unsure{\alpha|X}}^*$)

\[
E\simeq\underline{\operatorname{End}}(E)/\mathcal A\qquad(\simeq E^{\otimes_A2}\otimes_k L^\vee)
\]

équivaut à la \deleted{\unsure{iso}}donnée d'une section de

\[
E\otimes_k L^\vee,
\]

base de ce $\mathcal A$-module mod section provenant de $O^*$, Or il n'y en a pas de canonique, visiblement, \deleted{Donc} \inserted{i.e.} il n'y a pas de sous-fibré en droites canonique $D\subset E$ associée à la structure de similitude de départ sur $E$. Donc la supposition du début est erronée -- $\operatorname{End}(E)/\mathcal A$ est, modulo \deleted{\unsure{m}}Module inversible provenant de la base, un \deleted{\unsure{m}}Module inversible \uline{carré}, et c'est la racine carrée qu'il faudrait arriver à saisir \ldots
% archive sheet number: 31
\scanpage{33}

En rectifiant le tir, on trouve le $\mathcal A$-module

\[
E\otimes_{\mathcal A}(\unsure{{}_{\sigma}E|{}^{\sigma}E|\bar E|\check E})^{\otimes_A-1}\simeq\mathcal M/\deleted{\unsure{\mathcal A}}\,\mathcal M^{\natural}\quad\text{mod mod. inv. sur }\underline O
\]

où $\mathcal M^\natural$ est défini localement sur $X$ par le moyen d'une base $(1,u)$ de $\mathcal A$, via

\[
\underline{\varphi}_u:\mathcal M\mathbin{\unsure{-|\longrightarrow}}\mathcal M,\qquad\deleted{\unsure{x|\theta}}\unsure{\theta|\Theta}\longmapsto u\theta-\theta\bar u
\]

% slanted
\margin{\uline{NB} $\underline{\varphi}_u$ est $\mathcal A$-linéaire \unsure{p.|pour} les structures \unsure{NB} de $\mathcal A$-module \uline{gauche} sur $\mathcal M$}

($\underline{\varphi}_u$ dépend linéairement de $u$, et $\underline{\varphi}_1=0$, donc $\underline{\varphi}_u$ ne dépend que de la section \inserted{$\dot u$} de \deleted{\unsure{u}} $\Lambda^\vee\simeq\mathcal A/O$ définie par $u$ \ldots), en posant

\[
\mathcal M^\natural=\operatorname{Im}\underline{\varphi}_u,\quad\text{donc }\mathcal M/\mathcal M^\natural=\operatorname{Coker}\underline{\varphi}_u.
\]

De façon plus intrinsèque, on a suite exacte

\[
\deleted{\unsure{\mathcal A}}\mathcal M\otimes\Lambda^\vee\longrightarrow\mathcal M\deleted{\unsure{/\mathcal M^\natural}}\longrightarrow\Delta\longrightarrow0
\]

(\unsure{suite|visiblement} exacte de $A$-module, pour la structure de $A$-Module \uline{à gauche} de $\mathcal M$). Il faudrait ensuite définir sur $\Delta$ une forme \inserted{\unsure{bin.}} quadratique canonique, plus précisément un iso.

\[
N_{\mathcal A/\underline O}(\Delta)\simeq\underline O
\]

\unsure{Mettons|Notons} sur $\mathcal M$ l'autodualité donnée par

\[
\langle\theta,\theta'\rangle\overset{\mathrm{déf\unsure{n}}}{=}\operatorname{Tr}\theta\overline{\theta'}=\operatorname{Tr}\theta'\bar\theta
\]

alors, si $\lambda_{\unsure{v|u}},\gamma_v$ sont les opérations de la multiplication par $v$ à g. et à dr.

\[
\lambda_v\theta=v\cdot\theta,\quad \gamma_v\theta=\theta v,
\]

on a

\[
{}^t\lambda_v=\lambda_{\bar v},\quad {}^t\deleted{\unsure{\lambda}}\gamma_v=\gamma_{\bar v}
\]
% archive sheet number: 32
\scanpage{34}

donc la transposée de \deleted{\unsure{$\underline{\varphi}$}}

\[
\underline{\varphi}_u:\unsure{\theta|\Theta}\longmapsto u\theta-\theta\bar u
\]

est

\[
\deleted{\unsure{-|=}}\underline{\varphi}_{\bar u}:\theta\mathbin{\unsure{\mapstochar\relbar|\mapsto}}\bar u\theta-\theta u,
\]

or on a $\underline{\varphi}_{\bar u}=-\underline{\varphi}_u$ car $\bar u\equiv-u$ mod $\underline O\cdot1$ \inserted{et $\underline{\varphi}_1=0$}

Ceci dit, on voit donc que le dual \uline{sur $k$} de $\operatorname{Coker}\underline{\varphi}_u$ s'identifie à $\operatorname{Ker}\underbrace{\underline{\varphi}_{\bar u}}_{-\underline{\varphi}_u}$ ($=\operatorname{Ker}\underline{\varphi}_u$). Donc modulo \unsure{torsade|module|torsion|Module} inv. sur $\underline O$, $\operatorname{Coker}\underline{\varphi}_u$ s'identifie, en \deleted{tant} que $\mathcal A$-Module, au dual \inserted{sur $\mathcal A$} de $\operatorname{Ker}\underline{\varphi}_u$.

Soit

\[
\begin{aligned}
&\mathcal B\subset\mathcal M\\[-2pt]
&\ \Vert\\[-2pt]
&\left\{\theta\in\Gamma\mathcal M\ \middle|\ \bar u\theta-\theta u=0\quad
\substack{\text{pour une section}\\\text{locale de }\mathcal A\text{ qui}\\\text{forme base de }\mathcal A\\\text{avec }1,\text{ donc pour}\\\text{\uline{tte} section de }\mathcal A}\right\}
\end{aligned}
\]

On a

\[
\mathcal B\cap\mathcal M^*=\left\{\theta\in\Gamma\mathcal M^*\ \middle|\
\substack{\theta\text{ normalise }\mathcal A^*\text{ (ou }\mathcal A\text{),}\\\text{et induit sur }\mathcal A\text{ l'}\inserted{\text{i.e.}}\strut\\\text{invol. can. }\sigma}\right\}
\]

On voit donc que $\mathcal B\cap\mathcal M^*$ est un \inserted{pseudo}torseur sous $\underset{\substack{\Vert\\\mathcal A^*}}{\operatorname{Cent}_{\mathcal M^*}(\mathcal A)}\inserted{=\mathcal A^*}$,

Et c'est un torseur ssi $\exists$ section loc fppf. Mais \deleted{comme} OPS $\mathcal M=\mathbb M_2(O)$, et \deleted{comme} $\mathcal A=O[u]$, avec $u^2+bu+c=0$, et comme $\operatorname{Tr}\bar u=\operatorname{Tr}u=b\qquad N(\bar u)=N(u)=\unsure{c|2}$, on sait que $\exists\ \theta\in\mathcal M^*$ tel que $\deleted{\unsure{\theta}}\operatorname{Int}(\theta)\cdot u=\bar u$.
% archive sheet number: 33
